\section{Metodologia}
\label{sec:metodologia}

A presente pesquisa caracteriza-se como um estudo aplicado de natureza mista, estruturado a partir da integração de uma Revisão de Escopo (\textit{Scoping Review}) com o método de \textit{Design Science Research} (DSR) para o desenvolvimento de um artefato tecnológico. 

O fluxo metodológico inicia-se com uma etapa iterativa de busca e filtragem sistemática da literatura acadêmica, com o intuito de mapear a fundamentação arquitetural do \textit{Blockchain} e a sua quebra de paradigma na área de sistemas distribuídos. Esta revisão ocorre em um ciclo contínuo de três etapas: (1) busca de artigos base; (2) filtragem de relevância focada em inovações de redes descentralizadas; e (3) escrita e síntese do conhecimento. Uma vez consolidados os princípios de consenso em redes abertas, o ciclo avança para uma análise de domínio, na qual o escopo de busca é direcionado para estudos de caso focados especificamente nas indústrias financeira e da saúde. 

Por fim, após a finalização da escrita base e do mapeamento contextual de como o \textit{Blockchain} evoluiu para suprir as demandas operacionais destes setores, os resultados teóricos atuam como base de requisitos para o desenvolvimento prático. Fechando o fluxo metodológico via DSR, constrói-se uma plataforma computacional para a visualização de interações de rede \textit{peer-to-peer}, traduzindo os conceitos de descentralização e integridade mapeados na literatura em um ambiente de simulação e validação visual. A Figura \ref{fig:fluxo_metodologico} ilustra o pipeline da metodologia adotada.

\begin{figure}[!hbtp]
    \centering
    \resizebox{0.9\columnwidth}{!}{%
    \begin{tikzpicture}[
        % Estilos herdados do seu padrão
        recntagle/.style={
            rectangle,
            rounded corners=3pt,
            inner sep=6pt,
            text centered,
            font=\small\bfseries,
            draw=black!70,
            fill=blue!10,
            line width=0.8pt
        },
        title_block/.style={
            rectangle,
            rounded corners=3pt,
            inner sep=4.5pt,
            text centered,
            font=\small\bfseries,
            draw=black!70,
            fill=blue!10,
            line width=0.8pt
        },
        arrow/.style = {
            thick, 
            ->, 
            >={Stealth[length=2.5mm]}
        },
        box_delimitador/.style = {
            rectangle,
            rounded corners=8pt,
            draw=black!70,
            fill=black!2,
            line width=0.8pt
        },
        sub_block/.style = {
            rectangle,
            rounded corners=3pt,
            inner sep=5pt,
            text centered,
            draw=black!70,
            fill=blue!5,
            line width=0.7pt,
            font=\small\bfseries,
            text width=2.5cm,
            minimum height=1.2cm
        }
    ]

        % ==========================================
        % FASES 1 E 2: Revisão e Análise de Domínio
        % ==========================================
        \node (fases12_box) [box_delimitador, minimum width=10.5cm, minimum height=3.5cm] {};
        \node (fases12_title) [title_block] at (fases12_box.north) {Fases 1 e 2: Revisão de Escopo e Estudo de Domínio};

        % Nós internos do ciclo
        \node (busca) [sub_block] at ([xshift=-3.3cm, yshift=-0.2cm]fases12_box.center) {1. Busca de Artigos};
        \node (filtro) [sub_block] at ([xshift=0cm, yshift=-0.2cm]fases12_box.center) {2. Filtragem de Relevância};
        \node (escrita) [sub_block] at ([xshift=3.3cm, yshift=-0.2cm]fases12_box.center) {3. Escrita e Síntese};

        % Setas do fluxo principal
        \draw [arrow] (busca.east) -- (filtro.west);
        \draw [arrow] (filtro.east) -- (escrita.west);

        % Seta de retorno do ciclo iterativo
        \draw [arrow] (escrita.south) -- ++(0,-0.6cm) -| node[above, pos=0.25, font=\footnotesize\itshape] {Ciclo Iterativo} (busca.south);

        % ==========================================
        % FASE 3: Desenvolvimento Prático (DSR)
        % ==========================================
        \node (fase3) [recntagle, minimum width=8cm, below=1.2cm of fases12_box] {Fase 3: Desenvolvimento Prático da Plataforma (DSR)};

        % Conexão da Fase teórica para a Fase prática
        \draw [arrow] (fases12_box.south) -- (fase3.north);

    \end{tikzpicture}%
    }
    \caption{Fluxo metodológico da pesquisa, evidenciando o ciclo iterativo de revisão bibliográfica que fundamenta o posterior desenvolvimento prático da plataforma.}
    \label{fig:fluxo_metodologico}
\end{figure}
